\begin{enumerate}
  \item Linsmakarformeln ger
  \[
    P = (n - 1)\left(\frac{1}{r_1} - \frac{1}{r_2}\right) = 0{,}50\left(\frac{1}{0{,}20} + \frac{1}{0{,}30}\right) \approx 4{,}17\,\mathrm{D},
  \]
  så $f = 1/P = 24\,\mathrm{cm}$.
  \item Första ytan (vatten till glas, $r_1 = +0{,}20\,\mathrm{m}$) och andra ytan (glas till vatten, $r_2 = -0{,}30\,\mathrm{m}$):
  \[
    P_1 = \frac{1{,}50 - 1{,}33}{0{,}20} = 0{,}85\,\mathrm{D}, \qquad P_2 = \frac{1{,}33 - 1{,}50}{-0{,}30} \approx 0{,}57\,\mathrm{D},
  \]
  så $P = P_1 + P_2 \approx 1{,}42\,\mathrm{D}$, jämfört med $4{,}17\,\mathrm{D}$ i luft.
  \item Om man summerar ytornas ekvationer och använder $s_2 = -s_1'$ fås
  \[
    \frac{1{,}33}{s} + \frac{1{,}33}{s'} = P.
  \]
  Med $s \to \infty$ blir $s' = 1{,}33/1{,}42 \approx 0{,}94\,\mathrm{m}$, nästan fyra gånger längre än $24\,\mathrm{cm}$ i luft.
  \item Brytningen beror på \emph{skillnaden} i brytningsindex. Hornhinnan står för det mesta av ögats brytningsstyrka, och den har ungefär samma brytningsindex som vatten. Under vatten förlorar ögat därför nästan all sin brytning, och bilden hamnar långt bakom näthinnan. Simglasögon återställer ett luftlager framför ögat.
\end{enumerate}
